\subsection{Disposición en bloques y mapeo del vector de estado}

La partición del vector de estado se materializó mediante la clase
\texttt{ChunkLayout}, encargada de establecer la correspondencia entre índices lógicos
del arreglo de amplitudes y bloques físicos del backend comprimido. En la
implementación, esta clase calcula explícitamente: número total de estados,
número total de elementos reales almacenados (\texttt{2 * totalStates}), cantidad de
estados por bloque, elementos por bloque y número de bloques requeridos. Además
expone tres consultas decisivas: \texttt{chunkElemCount}, para conocer el tamaño real de
un bloque; \texttt{chunkStateBegin}, para ubicar el primer estado lógico de cada bloque; y
\texttt{stateRef}, para traducir un índice global a \texttt{(chunkIndex, elemOffset)}.

Este mapeo cumple una función central dentro de la implementación. Por una parte,
encapsula la granularidad del almacenamiento y evita que otras partes del sistema
reconstruyan manualmente la relación entre estados y bloques. Por otra, garantiza un
tratamiento correcto del último bloque, cuya cantidad efectiva de amplitudes puede ser
menor que la nominal cuando $2^n$ no es múltiplo exacto de \texttt{chunkStates}. Esto es
relevante porque el backend opera sobre amplitudes complejas intercaladas
\texttt{[real, imag, real, imag, \ldots]}, y por tanto cada estado ocupa exactamente dos
elementos dentro del bloque.

En términos operativos, el backend no trabaja con el vector completo como un único
buffer comprimido, sino con una secuencia indexada de fragmentos. Cada consulta o
modificación pasa primero por \texttt{ChunkLayout}: si se solicita la amplitud de un estado,
la clase devuelve el bloque que debe cargarse y el desplazamiento local del par
real/imaginario; si se recorre el vector por bloques, también permite conocer cuántos
elementos deben leerse o escribirse en cada fragmento. Gracias a esta mediación, la
recuperación localizada y la escritura posterior pueden resolverse sin materializar el
vector completo en memoria densa.
